% Part: model-theory
% Chapter: lindstrom
% Section: ls-property

\documentclass[../../../include/open-logic-section]{subfiles}

\begin{document}

\olfileid[fr]{mod}{lin}{lsp}

\olsection{Propriétés de compacité et de Löwenheim--Skolem}

Étendons aux logiques abstraites les
propriétés de compacité et de
Löwenheim--Skolem.

\begin{defn}
Une logique abstraite
$\tuple{L,\models_L}$ a la
\emph{propriété de compacité} si,
pour tout langage~$\Lang{L}$ et tout
ensemble $\Gamma\subseteq L(\Lang{L})$,
la satisfaisabilité de chaque
sous-ensemble fini de~$\Gamma$
entraîne celle de~$\Gamma$.
\end{defn}

\begin{defn}
$\tuple{L,\models_L}$ a la
\emph{propriété de Löwenheim--Skolem
descendante dénombrable} si, dans tout
langage au plus dénombrable, chaque
ensemble au plus dénombrable et
satisfaisable d'énoncés de
$L(\Lang{L})$ possède un modèle
au plus dénombrable.\footnote{Précision éditoriale : la source demande un modèle dénombrable pour tout ensemble satisfaisable, sans restriction de langage. Cette propriété est ici formulée pour les théories dénombrables en langage dénombrable, comme le requiert le codage utilisé. Elle porte sur une théorie entière, et non sur un seul énoncé ; certaines présentations nomment cette hypothèse « propriété de Skolem dénombrable ».}
\end{defn}

La notion d'isomorphisme partiel de la
\olref[bas][pis]{defn:partialisom}
est algébrique : elle porte sur les
interprétations des symboles, non
sur la forme des énoncés. Elle
s'applique donc aux structures
des logiques abstraites. Pour les
structures relationnelles sans
constantes, le
\olref[bas][pis]{thm:p-isom2}
montre que l'isomorphisme partiel
entraîne l'équivalence élémentaire
au premier ordre; le même raisonnement
tient compte des constantes en
préservant leurs interprétations.
Cette récurrence sur les !!{formula}s
du premier ordre ne s'étend pas
automatiquement à un énoncé
arbitraire~$!E\in L(\Lang{L})$.
La propriété de Löwenheim--Skolem
permet néanmoins de conclure.

\begin{thm}
\ollabel{thm:abstract-p-isom}
Soit $\tuple{L,\models_L}$ une
logique normale ayant la propriété
de Löwenheim--Skolem descendante
dénombrable. Deux !!{structure}s
partiellement isomorphes sont alors
élémentairement équivalentes dans
$\tuple{L,\models_L}$.
\end{thm}

\begin{proof}
Supposons $\Struct{M}\simeq_p\Struct{N}$,
mais $\Struct{M}\models_L !E$ et
$\Struct{N}\not\models_L !E$ pour
un énoncé~$!E$. Par la propriété
de dépendance finie, on peut
remplacer~$!E$ par un représentant
dans un langage relationnel
\emph{fini} et réduire les deux
structures à ce langage. Les
réduits restent partiellement
isomorphes. L'invariance par
isomorphisme permet en outre
d'en prendre des copies de
domaines disjoints.\footnote{Compléments éditoriaux : la relativisation de $!E$ produit des énoncés de la logique abstraite, que la source appelle à tort des énoncés du premier ordre. Nous développons le codage du va-et-vient en une théorie dénombrable, avec une relation par longueur de tuple, au lieu de la construction condensée par suites de la source. La figure est adaptée à ce codage.}

Formons un langage auxiliaire
au plus dénombrable. Il contient
deux copies des symboles du
langage fini, deux !!{predicate}s
unaires $U,V$ et, pour chaque
$k\geq1$, un !!{predicate}~$I_k$
à $2k$ places. Dans une structure
de ce langage, $U$ et~$V$
découpent deux sous-structures
$\Struct{M}_U$ et~$\Struct{N}_V$;
les constantes de chaque copie
appartiennent au côté correspondant.
Par renommage et relativisation,
on obtient des énoncés abstraits
$!E_U$ et~$!E_V$ qui expriment
respectivement que
$\Struct{M}_U\models_L !E$ et
$\Struct{N}_V\models_L !E$.

Ajoutons un ensemble dénombrable
$\Theta$ d'énoncés \emph{du premier
ordre}, représentés dans la logique
normale par les équivalents de la
remarque précédente. Ils décrivent
une famille de
va-et-vient. Cet ensemble impose que $U,V$
soient disjoints et non vides.
Pour chaque~$k$, $I_k(\mathbf a,
\mathbf b)$ ne vaut que pour des
$k$-uplets des deux côtés et
entraîne qu'ils définissent un
isomorphisme partiel : mêmes
égalités, mêmes rapports aux
constantes et mêmes formules
atomiques du langage fini.
L'instance vide impose, pour
chaque $a\in U$, un $b\in V$
avec $I_1(a,b)$, et réciproquement.
Pour tout $k\geq1$, les clauses
de va-et-vient imposent
\[
\begin{aligned}
I_k(\mathbf a,\mathbf b)
&\Longrightarrow
  (\forall a\in U)(\exists b\in V)
  I_{k+1}(\mathbf a a,\mathbf b b),\\
I_k(\mathbf a,\mathbf b)
&\Longrightarrow
  (\forall b\in V)(\exists a\in U)
  I_{k+1}(\mathbf a a,\mathbf b b).
\end{aligned}
\]
Comme le langage de départ est
fini, chaque condition atomique
pour une longueur~$k$ se formule
par un nombre fini d'axiomes
du premier ordre; l'ensemble de
toutes ces clauses est dénombrable.

\begin{figure}[h]
  \centering
  \begin{tikzpicture}[node distance=2cm, auto, thick, >=stealth']
    \draw [rounded corners] (0,0) -- (8,0) -- (8,4) -- (0,4) -- cycle;
    \draw (2,2) circle (1.1cm);
    \draw (6,2) circle (1.1cm);
    \path node at (0.75,3.5) {\large $\Struct{C}$};
    \path node at (2,2) {\large $\Struct{M}_U$};
    \path node at (6,2) {\large $\Struct{N}_V$};
    \node (Idom) at (3.1,2) {};
    \node (Irng) at (4.9,2) {};
    \draw[<->, bend left] (Idom) to node {\large $I_k$} (Irng);
  \end{tikzpicture}
  \caption{Dans la structure auxiliaire $\Struct{C}$,
    les relations $I_k$ codent le va-et-vient entre
    les deux sous-structures.}
\end{figure}

La théorie dénombrable
$\Theta\cup\{!E_U,\lnot !E_V\}$
est satisfaisable. En effet,
prenons la réunion disjointe des
copies de~$\Struct{M}$ et
$\Struct{N}$; interprétons les
$I_k$ par les couples de $k$-uplets, avec
répétitions permises, dont la correspondance est une
restriction d'un membre de la famille de va-et-vient
contenant les valeurs de toutes les constantes.
Ces restrictions forment encore une famille
de va-et-vient et incluent la fonction vide.
En effet, on peut prolonger tout membre
pour lui adjoindre les valeurs des constantes,
en nombre fini, puis le prolonger encore
pour traiter un élément supplémentaire.
Les restrictions ainsi obtenues préservent aussi
les formules atomiques contenant ces constantes. Les conditions de~$\Theta$
et les deux énoncés relativisés
sont vérifiés.

Par Löwenheim--Skolem, cette
théorie a un modèle au plus
dénombrable~$\Struct{C}_0$.
Ses côtés $\Struct{M}_0$ et
$\Struct{N}_0$ restent séparés
par~$!E$, tandis que les
relations $I_k$ y fournissent
une famille non vide
d'isomorphismes partiels avec
va-et-vient. Ces côtés sont donc
partiellement isomorphes et au
plus dénombrables. Le
\olref[bas][pis]{thm:p-isom1}
les rend isomorphes, en
contradiction avec l'invariance
de~$!E$ par isomorphisme.
\end{proof}

\end{document}
